% ---------------------------------------------------------------------
%  Macros de notación (la del profesor).
%  Si cambias una notación aquí, cambia en todo el documento.
% ---------------------------------------------------------------------

% Orden de error: \Order{h^2}  ->  O(h^2)
\newcommand{\Order}[1]{\mathcal{O}\!\left(#1\right)}

% Diferencial recto y derivada temporal: \ddt{u} -> du/dt
\newcommand{\dd}{\mathrm{d}}
\newcommand{\ddt}[1]{\frac{\dd #1}{\dd t}}

% Valor EXACTO en un nodo:   \uex{\ell}  -> u(t_\ell)
% Valor NUMÉRICO en un nodo: u_\ell      (se escribe directo)
\newcommand{\uex}[1]{u(t_{#1})}

% Punto medio entre nodos: t_{\lmed} -> t_{\ell+1/2}
\newcommand{\lmed}{\ell+1/2}
% Igual con índice n (notación del examen): t_{\nmed} -> t_{n+1/2}
\newcommand{\nmed}{n+1/2}

% Punto fijo / solución particular
\newcommand{\ufijo}{u^{*}}

% Orden de convergencia global / local (Sesión 02)
\newcommand{\OCG}{\mathrm{OCG}}
\newcommand{\OCL}{\mathrm{OCL}}

% Matrices y vectores (para las Partes II–IV)
\newcommand{\mat}[1]{\mathbf{#1}}
\newcommand{\vect}[1]{\bm{#1}}
\newcommand{\tridiag}{\operatorname{tridiag}}
\newcommand{\keff}{k_{\mathrm{ef}}}
